% ============================ 1. INTRODUCTION ============================
\section{Introduction}

\subsection{A city that rides, and the problem that follows}

Between 2016 and 2018, dockless shared bicycles changed how Chinese cities moved.
Unlike the earlier station-based systems of Hangzhou and Shanghai, a dockless fleet
imposes no docking constraint: a rider may leave the bicycle wherever the journey
ends, and the next rider finds it through a phone application. The convenience is
real, and so is the externality. Because no station absorbs the end of a trip, the
spatial distribution of bicycles is simply the integral of every rider's private
choice, and that distribution drifts. Bicycles pile up at subway exits, at shopping
centres and at the foot of office towers, while residential streets empty by nine in
the morning. Operators respond with trucks; cities respond with parking rules; and
both are expensive because neither is guided by a model of where the bicycles will
go.

This paper builds that model from the largest artefact the industry left behind: the
2017 Mobike Algorithm Challenge dataset, which records 3,213,310 orders placed on
484,754 bicycles by 349,241 users in Beijing between 10 and 24 May 2017. Each order
carries an order identifier, a user identifier, a bicycle identifier, a start time,
a bicycle type, and origin and destination coordinates encoded as seven-character
geohashes. There is no inventory record, no station, and no record of a user who
wanted to ride and found no bicycle.

\subsection{The five questions}

The problem poses five questions, which we restate in operational terms:

\begin{enumerate}[label=\textbf{Q\arabic*.},leftmargin=2.2em]
\item \textbf{Behaviour.} Build a mathematical model of how users ride: when they
  start, where they go, and how long it takes.
\item \textbf{Rebalancing.} Design an optimal repositioning and scheduling plan for
  a typical day, using trucks or tricycles to move bicycles between areas.
\item \textbf{Incentives.} Explore a scheme that pays users to end their trips where
  the platform wants bicycles.
\item \textbf{E-fences.} Predict parking demand by location and propose a layout and
  capacity plan for electronic fences.
\item \textbf{Competition.} Discuss how the coexistence of competing platforms
  affects the resolution of the problems above.
\end{enumerate}

\subsection{Three things the data do not tell us}

Most of the analytical difficulty in this problem is not in the optimisation; it is
in three quantities the dataset does not contain and that a careless model would
silently assume.

\textbf{No inventory.} The data record movements, not stocks. A cell that receives no
arrivals may be empty or may hold fifty bicycles that nobody rode. Any statement
about a ``shortage'' therefore requires an assumption about where the bicycles were
at midnight, and that assumption is not identified by the data (Section~\ref{sec:sens}).

\textbf{Unknown sampling rate.} The file is described as a sample of the operator's
orders. If only a fraction $p$ of orders is present, then every count in the file is
a count of $1/p$ of reality, and the fleet is larger than the number of distinct
bicycle identifiers. We estimate $p$ from the internal consistency of successive
rides of the same bicycle, and we show that the fleet estimate is invariant to it.

\textbf{Truncated demand.} A user who wants to ride and finds no bicycle never
appears. The observed trips are therefore demand censored at supply, which biases any
naive estimate of demand downward exactly where the system is most stressed. Our
models avoid this bias by treating observed trips as realisations of a rate process
rather than as the demand itself, and by reporting service quality as a simulated
quantity rather than a fitted one.

\subsection{Related work}

Bike-share research divides into three strands. The first documents the systems and
their users \cite{shaheen2010,fishman2013}; studies of station-based and free-floating
systems add land-use effects on bicycle flows \cite{faghih2014}, predictors of
shared-bicycle use \cite{bachand2012}, and dockless travel patterns in Singapore and
Nanjing \cite{shen2018,du2018}. The second treats rebalancing as a vehicle-routing or
inventory problem: static repositioning \cite{raviv2013}, the static rebalancing
problem with multiple vehicles \cite{chemia2013}, inventory-plus-routing formulations
with service-level constraints \cite{schuijbroek2017}, and stationless static
rebalancing for free-floating systems \cite{pal2017}. The third replaces trucks with
prices, either by dynamic redistribution with user incentives \cite{pfrommer2014} or
by mechanism design for balanced systems \cite{singla2015}. The probabilistic
analysis of station occupancy that underlies both is given in \cite{omahony2015}.

Our approach differs from all three in one respect: the behavioural model is not a
convenience for generating demand but the object of study in its own right, and the
same estimated object is then reused, without re-fitting, in the optimisation (Q2),
the incentive design (Q3) and the parking forecast (Q4). Spatial interaction models
of the gravity family \cite{wilson1970,erlander1990,fotheringham1989} supply the
destination choice; the discrete-choice foundation is due to \cite{mcfadden1974}.
The radiation model offers an opportunity-based alternative \cite{simini2012}. The
latent-class segmentation uses $k$-means \cite{lloyd1982} with the silhouette criterion
\cite{rousseeuw1987}. The competition discussion draws on the economics of spatial
and platform competition \cite{hotelling1929,katz1985,rochet2003}.

\subsection{Executive summary: contributions and answers}

The contributions are four.

\begin{enumerate}[leftmargin=1.6em]
\item \textbf{A diagnosis.} The Beijing system of May 2017 was oversupplied. We
  establish this three ways --- a capture--recapture fleet estimate, a zone-level
  stockout diagnostic, and a cell-level simulation --- and it reframes Q2 from
  ``prevent shortages'' to ``minimise the cost of restoring balance''.
\item \textbf{One model, four uses.} A single estimated gravity model feeds the
  rebalancing LP, the incentive supply curve and the parking forecast, so the four
  answers are internally consistent.
\item \textbf{An incentive supply curve.} Rather than assert an acceptance rate, we
  derive the cost of voluntary relocation as a function of the reward, solve an hourly
  assignment program with one reroute per trip, and price the instrument against both
  layers: it cannot replace the trucks, and it can replace part of the tricycle layer.
\item \textbf{Explicit uncertainty.} The unidentified quantities are bounded, not
  assumed: the sampling rate is estimated and its role as an upper-bound scale factor
  is stated, the initial-inventory proxy is carried as two scenarios, and every behavioural
  parameter is swept in Section~\ref{sec:sens}.
\end{enumerate}

\tabref{tab:answers} states the answer to each question in one line.

\begin{table}[htbp]
\centering
\footnotesize
\caption[Answers at a glance]{Answers at a glance. All figures are for the typical weekday defined in
Section~\ref{sec:data}; costs are in 2017 CNY.}
\label{tab:answers}
\footnotesize
\begin{tabular}{@{}p{0.7cm}p{4.1cm}p{5.0cm}p{4.6cm}@{}}
\toprule
 & \textbf{Question} & \textbf{Model} & \textbf{Answer} \\
\midrule
Q1 & Behaviour & NHPP arrivals + singly-constrained gravity + speed-based duration & Median trip 0.70 km; distance dominates (log-loss 3.24--3.52 against 8.57 uniform); one shared $\beta$ is an acceptable simplification \\
Q2 & Rebalancing & Integer time-space network MILP (trucks) + same-zone tricycle routing & Two-tier: 60 trucks + 975 tricycles, 617{,}055.7515 CNY/day, saving 902{,}223.20 (59.385\%), unmet 39{,}984.2332 \\
Q3 & Incentives & Logit acceptance + hourly assignment program & Cannot replace trucks (at least 1{,}532 CNY/day worse at zero payment); can replace tricycles: $r^*=0.60$, saving 18{,}011.46 (2.92\%) \\
Q4 & E-fences & Greedy maximal coverage + Monte-Carlo peak occupancy & 4{,}000 sites $\times$ 150 slots: 59.53\% occupancy / 55.80\% arrival coverage at 343{,}638; 394{,}553.13 unassigned \\
Q5 & Competition & Two-platform cell-resolution simulation & Unmet 12.10\% against 11.66\% single; capacity 1.08--1.13$\times$; sharing does not reduce lost orders \\
\bottomrule
\end{tabular}
\end{table}

\subsection{Roadmap}

Section~\ref{sec:assume} states assumptions and notation.
Section~\ref{sec:data} describes the data, the geohash decoding, the sampling-rate
estimate and the fleet-size estimate.
Section~\ref{sec:q1} builds and validates the behavioural model.
Sections~\ref{sec:q2}--\ref{sec:q5} answer Q2--Q5.
Section~\ref{sec:sens} reports sensitivity and the boundaries of validity, and
Section~\ref{sec:eval} evaluates the work.
